% 第 3 章　图像会说话
\chapter{图像会说话}

\step{反常识开场}

先看一场百米决赛的两种记录。

第一种记录是冲线照片：选手甲的胸口先压过终点线，领先乙大约半个身位。裁判宣布，甲是冠军。没有人有异议。

第二种记录是测速雷达：它在整场比赛中每隔零点几秒测一次两人的速度。把数据摊开看，你会发现一件怪事——从发令枪响后第三秒开始，到冲线前一刻为止，乙的速度几乎一直比甲高。乙起跑慢了，但中途追得很凶，最高速度也比甲高。

那么问题来了：这场比赛里，到底谁“跑得更快”？

如果“更快”的意思是“谁先跑完一百米”，答案是甲。如果“更快”的意思是“谁的速度更大”，那么在大半段时间里答案是乙。两种说法都没错，却给出不同的冠军。日常生活中我们随口说“他跑得比我快”，其实在这两种意思之间滑来滑去，从来没有分清过。

物理学家遇到这种“一句话两种意思”的局面时，不会靠吵架解决，而是换一套更精确的语言。本章要介绍的这套语言你早就见过：画一张图。你会看到，只要把这场比赛画成图，甲和乙的区别会一眼显现出来，根本不需要争论；你还会看到，图上一根线的“陡”和图上一块“面积”，分别对应着物理世界里最重要的两类问题——“现在变得有多快”和“一共积累了多少”。这两个问题将贯穿整本书：从运动到力，从电流到热量，最后一直到量子力学。

顺便说一句：这套语言的发明本身是一场革命。在它出现之前，人们只能用文字描述变化——“越来越快”“渐渐变慢”——谁也说不清“越来越快”到底有多快。在它出现之后，变化第一次变成了可以精确谈论、可以计算、可以画出来的东西。

\step{先猜一猜}

在往下看之前，先把你自己的猜想写下来。本章结束时你可以回来对照，看看直觉对了多少。

第一题。一只透明杯子放在水龙头下接水，水流大小保持不变。杯里的水面高度随时间怎么变化？如果画一张“水面高度—时间”图，你猜它是直线、向上弯的曲线，还是向下弯的曲线？

第二题。电梯从一楼出发，中途在四楼停了一次，最后到八楼。如果画一张“电梯位置—时间”图，中途停留的那十几秒，图上的线是什么样子？是继续往上走，还是拐个弯？

第三题。汽车的仪表盘上有两块表：时速表显示此刻的速度，里程表显示一共开了多远。假设你踩下油门，时速表指针从每小时四十公里爬到八十公里。这段时间里，里程表增加的数字，跟时速表指针扫过的“快慢”有什么关系？能不能从速度的记录里算出里程，反过来，能不能从里程的记录里算出速度？

第四题，也是最微妙的一题。一个人先向北走，再掉头向南走回原处。如果画“速度—时间”图，他向南走的那一段，图线应该画在横轴上方还是下方？你的直觉说“速度怎么会有负的”——先记住这个反应，本章后面它会被正式推翻。

不必急着回答。接下来我们做一个实验，让你的手和眼先于你的头脑得到答案。

\step{动手实验}

\begin{experiment}[滴漏计时：亲手画出一条直线]
你需要：一只透明玻璃杯（或剪开的矿泉水瓶）、一支记号笔、一把尺子、一部手机（用它的秒表）、一个能滴水的龙头或扎了小孔的塑料杯。

第一步，把水龙头调到很小但稳定的水流，让水匀速流进玻璃杯。

第二步，水流稳定后启动秒表。每隔二十秒，用记号笔在杯壁上标出当时的水面位置，一共标六到八次。动作要快，标完继续计时，不要停表。

第三步，关掉水，用尺子量出每个标记到杯底的高度，再记下每个标记对应的秒表读数。你得到了一张两列数字的表：时间，以及对应的水面高度。

第四步，找一张方格纸（或自己画格子），横轴画时间，纵轴画水面高度，把每一对数字点成一个点，然后看看这些点的排列。

如果水流足够稳定，你会看到这些点几乎排成一条直线。这是本章最重要的画面之一：\textbf{均匀的变化画出来是一条直线}。直线越陡，说明水面升得越快，也就是水流越大。

再做一个对比：用手捏住软管改变水流，先大后小，重复实验。这次点会排成一条先陡后平的折线。你已经无师自通地画出了本章的主角——一张会说话的图像。
\end{experiment}

这个实验还有一个隐藏版本，留到本章末尾的挑战题里：如果把直桶玻璃杯换成一个上粗下细的杯子，水面高度还是直线上升吗？先别翻答案。

\step{旧模型遇到麻烦}

你可能会问：我们已经有文字，有表格，还有“平均速度”这个初中就学过的概念，为什么还嫌不够？因为这三样东西各自有一个治不好的毛病。

先说文字。“列车出站后越来越快，快到站时又渐渐慢下来。”这句话人人能懂，但它回答不了任何具体的问题：快了多少？用了几秒变快的？哪一段变得最急？文字描述变化，就像只用“挺高的”“相当高”来描述人的身高——日常寒暄够用，订做西装就不行了。

再说表格。表格确实精确，但它太诚实了，什么都告诉你，等于什么都没告诉你。一段高铁运行记录每秒钟一个速度值，十分钟就是六百个数字。盯着六百个数字，你很难“看见”列车是在第三分钟还是第七分钟开始减速的。规律藏在一堆数字里，人的眼睛不擅长从数字队列里看出形状。

最后说平均速度，这是最隐蔽的麻烦。一辆公交车全程十公里开了三十分钟，平均速度是每小时二十公里。这句话是真的，但它把一路上所有的故事都抹平了：等红灯时速度是零，刚出站时慢，过长江大桥时快。平均速度就像用“平均水深半米”来描述一条河——对游泳的人来说，这句话可能是致命的，因为河中间可能有三米深。平均值回答“总体上如何”，却故意不回答“此刻如何”。而物理恰恰经常需要知道“此刻”：撞车那一刻的速度决定安全气囊要不要弹出，落地那一刻的速度决定鸡蛋碎不碎。

更麻烦的是“此刻的速度”这个说法本身。速度是“一段时间内走了多远”除以“这段时间”，那“某一瞬间”——长度为零的一段时间——的速度是什么意思？零除以零吗？我们的旧概念在这里直接撞墙了。

你每天都在使用的导航软件，就是平均速度模型的重度用户。它预计“还有四十分钟到达”，靠的无非是把剩下的路程除以它估计的各路段平均速度。一旦前方堵成一团——车流速度忽快忽慢——这个预计就开始失灵，到达时间一跳一跳地往后改。不是导航变笨了，而是“用平均掩盖此刻”这种模型在剧烈变化面前天然会失手。它给你的其实是一句平均意义上的话，而你想知道的是每一刻的真实情况。

历史上，人们在这堵墙前卡了将近两千年。古希腊的芝诺用“飞箭不动”的悖论嘲笑过它：飞箭在每一个瞬间都占着一个确定的位置，那它凭什么在动？要拆掉这堵墙，需要的不是更聪明的诡辩，而是一个全新的概念装置。

\step{发明新概念}

这个装置分三层，我们一层一层把它造出来。

第一层：\textbf{坐标}。要谈论“在哪里”，先给每个位置起一个名字。沿着一条路竖起一根带刻度的轴，选一个点叫零，规定一个方向叫正，于是每个位置都得到一个数：零右边三米叫 \(+3\)，零左边两米叫 \(-2\)。注意负号在这里不表示“比没有还少”，它表示“在零点的另一侧”。这就是第四题里“负速度”的伏笔：正负是方向，不是多少。这一步看起来平淡，实则大胆——它宣布位置这种东西可以用数来记账。

还要强调一件容易被忽略的事：零点放在哪里、哪边算正，完全是你的自由。把跑道终点设为零点，把向西设为正方向，所有的数都会变，但图线的形状、斜率的符号关系、两块面积之差——也就是所有的物理内容——一概不变。坐标是我们贴在世界上的一张标签纸，不是世界本身的一部分。这个“物理规律不依赖标签贴法”的念头现在看着不起眼，到后面讨论参考系和相对论时，会长成全书最重的一根支柱。

第二层：\textbf{变量与函数}。在滴漏实验里，你记下的每一行都是一对数：一个时刻，一个高度。只要规则是确定的——“给我一个时刻，我告诉你那一刻的高度”——我们就说高度是时间的函数。函数是一台什么样的机器？它像一台自动售货机：投进一个数，必定吐出一个确定的数，不多不少，也不会看心情。这是它像的地方。它不像的地方也要说清：自动售货机里并没有真的“装着”高度，函数也不是物体肚子里藏着的一卷纸条，它只是我们用来记录和预言对应关系的一本账。另外，函数关系不一定是因果关系——“下午三点气温二十度”并不意味着三点钟造成了二十度，这一点后面讨论实验数据时还会回来。

第三层：\textbf{图像}。账本上每一对数，都可以点成纸上的一个点：横着数时间，纵着数高度。把所有点连起来，就是图像。图像的魔力在于，它把整本账压缩成一眼可见的形状。六百个高铁速度数字看不出名堂，六百个点连成线，列车在第几分钟减速一目了然。文字讲不清的、表格埋起来的，图像用形状直接喊出来。

形状里有两个密码，是本章真正的主角。

密码一：\textbf{斜率，回答“变得有多快”}。在图线上取两点，纵向差除以横向差，就是斜率。位置—时间图上，纵向差是位置的变化，横向差是时间的变化，斜率就是“单位时间内位置改变多少”——这正是速度。图线越陡，变化越快；图线水平，斜率是零，说明什么都没变；图线往下走，斜率是负，说明在往负方向变化。回想滴漏实验：水流恒定，水面均匀上升，图是直线，斜率处处相同；水流先大后小，图线先陡后平，斜率先大后小。斜率就是“陡”的精确版本。

斜率像什么？它像山坡的陡度：爬升的高度除以前进的水平距离。这是它像的地方。不像的地方有两个：其一，真实山坡的陡度受土质限制，而图像的斜率可以要多大有多大，图线竖起来时斜率趋于无穷——物理上那意味着“不花时间就完成了变化”，通常是模型出错的信号；其二，山坡的陡是眼睛看到的，图像的陡却依赖坐标轴的比例尺，同一组数据把横轴拉长就显得平缓。所以斜率必须带着单位报出来，比如“每秒三米”，光说“很陡”没有意义。

密码二：\textbf{面积，回答“积累了多少”}。换一个角度。如果知道的是每一刻的速度，想知道一共走了多远，怎么办？把时间切成许多小段，每一小段里速度差不多没变，位移大约等于“速度乘以这一小段时间”。在速度—时间图上，“速度乘以时间”恰好是一个细高矩形的面积：高是速度，宽是时间。把所有小矩形的面积加起来，就是图线与横轴之间的总面积。于是：\textbf{速度—时间图下的面积，就是这段时间内位置的总变化}。匀速行驶时，图线是一条水平线，面积是一个矩形，“速度乘时间等于路程”——你小学就会的公式，原来是面积的一个特例。

面积像什么？它像水表的读数：水表不直接告诉你此刻水流多大，它把每一刻的流量默默累加起来。这是它像的地方。不像的地方在于：图上的“面积”不是用尺子量出来的平方厘米，它的单位是纵轴单位乘以横轴单位——速度乘时间得到的是米，不是平方米。几何只是它借用的外形。

还有一个更漂亮的对应值得单独记住。斜率和面积是互为逆运算的一对：从位置图量斜率得到速度，从速度图量面积回到位置。汽车的仪表盘把这个对应摆在你眼前——时速表是一块“斜率表”，里程表是一块“面积表”，两块表永远互相吻合。这正是后面数学桥层要说的、整个微积分最核心的思想。

\step{一句话模型}

\begin{oneline}
位置—时间图上，线的斜率告诉你“此刻变得有多快”；速度—时间图下，面积告诉你“一共积累了多少”——斜率管“快慢”，面积管“总量”，两者互为逆运算。
\end{oneline}

用这句话去扫描世界，你会发现它无处不在。水表电表记录的是面积（积累），水流动静和电热水器的工作状态是斜率（快慢）；手机健康软件里，“此刻配速”是斜率，“今日总步数折算的距离”是面积；银行账户里，每月工资进出的速率是斜率，余额是面积。物理学的很多分支，说到底就是在不同场合反复使用这一对工具。

再提醒一句：这套语言并不专属“运动”。医院心电图的纵轴是电压、横轴是时间，医生读的其实是电压变化的形状；天气预报里那条气温曲线，斜率为正的上午在升温，斜率转负的傍晚在降温；电饭煲的功率—时间图下面积，就是这一次煮饭消耗的电能，也是你家电费单上“度”的来历。任何一个量只要随另一个量变化，图像语言就适用。本章用运动来讲它，只是因为运动最看得见摸得着。

\step{数学翻译器}

现在把上面的直觉一条一条翻译成符号。请记住顺序：先有画面，后有式子。

第一条翻译：斜率。在位置—时间图上取两个时刻 \(t_1\) 和 \(t_2\)，对应位置 \(x_1\) 和 \(x_2\)。这段时间里的平均速度就是图线上两点连线的斜率：
\[
\bar{v}=\frac{\Delta x}{\Delta t}=\frac{x_2-x_1}{t_2-t_1}.
\]
先按老规矩检查它。位置没变，\(\Delta x=0\)，平均速度是零——停着的人确实没有平均速度，对。时间越短而位移相同，分母越小，数值越大——同样的路程跑得越快，对。掉头往回走，\(x_2<x_1\)，\(\Delta x\) 为负，平均速度为负——符号忠实地报告了方向反转，先猜一猜里的第四题到此有了答案：向南走的那一段，速度—时间图线就应该画在横轴下方。

\begin{center}
\begin{tikzpicture}[>=Stealth,scale=1.05]
  \draw[->] (0,0) -- (5.6,0) node[below left] {\(t\)（时间）};
  \draw[->] (0,0) -- (0,3.4) node[below left] {\(x\)（位置）};
  \draw[thick] (0,0.4) -- (4.8,3.0);
  \draw[dashed] (1.5,1.19) -- (3.8,1.19) node[midway,below] {\(\Delta t\)};
  \draw[dashed] (3.8,1.19) -- (3.8,2.50) node[midway,right] {\(\Delta x\)};
  \fill (1.5,1.19) circle (1.2pt) (3.8,2.50) circle (1.2pt);
  \node[align=center] at (1.9,2.9) {斜率 \(\displaystyle=\frac{\Delta x}{\Delta t}\)};
\end{tikzpicture}
\end{center}

第二条翻译：从平均到瞬时。平均速度的麻烦在于它把一段时间抹平了。补救办法很直接：把这段时间缩短。测一秒钟内的平均速度，比测十秒的更接近“此刻”；测十分之一秒，又更接近。一路缩短下去，数值会稳定地逼近一个确定的数——这个数就定义为那一时刻的\textbf{瞬时速度}。直觉上，它就是图线在那一点处切线的斜率：想象把图线那一小段不断放大，曲线越来越像直线，这条直线的斜率就是瞬时速度。

注意这个定义做了什么：它没有去算“零除以零”，而是算一串越来越短的时间间隔里的平均速度，看这串数趋向哪里。芝诺的飞箭就是这样被赦免的——“瞬间的速度”不再是自相矛盾的说法，而是一个极限过程的缩写。这里没有任何玄学，只是约定：当间隔越取越小、读数稳定到不再变化时，我们就把这个稳定值叫作瞬时速度。测量的仪器当然永远只能取有限短的间隔，但这不妨碍定义本身干净成立，就像你画不出完美的圆，不妨碍“圆”这个概念是干净的。

第三条翻译：面积。知道每一刻的速度 \(v\)，求从 \(t_1\) 到 \(t_2\) 位置变化了多少。把时间切成许多小段，每段宽 \(\Delta t\)，段内速度近似不变，该段的位移近似是 \(v\,\Delta t\)，加起来：
\[
\Delta x\approx v_1\,\Delta t+v_2\,\Delta t+v_3\,\Delta t+\cdots
\]
小段切得越细，近似越好，极限情况下就是速度—时间图线下的面积。

\begin{center}
\begin{tikzpicture}[>=Stealth,scale=1.05]
  \draw[->] (0,0) -- (5.6,0) node[below left] {\(t\)（时间）};
  \draw[->] (0,0) -- (0,3.4) node[below left] {\(v\)（速度）};
  \fill[blue!15] (0,0) -- (0,0.8) -- (4.2,2.9) -- (4.2,0) -- cycle;
  \draw[thick] (0,0.8) -- (4.2,2.9);
  \draw[dashed] (4.2,2.9) -- (4.2,0);
  \node[below] at (4.2,0) {\(t_2\)};
  \node at (1.9,1.05) {面积 \(=\) 位置的变化};
\end{tikzpicture}
\end{center}

照例检查。速度为零，图线贴着横轴，面积为零，位置不变——停车时里程表不走，对。速度恒定为 \(v\)，图线是水平线，面积是矩形 \(v(t_2-t_1)\)，回到“速度乘时间”的小学公式，对。速度为负时，图线在横轴下方，按规定面积记为负，位置的变化也是负——车倒回去，里程表的“位移账”在减少，对。注意这里说的是位移账，不是车坏了的里程表；这个细节本章末尾还会回来。

第四条翻译是前两条的合奏，也是一套真实的工程工具。铁路部门调度列车用的“列车运行图”，本质上就是一张位置—时间图：横轴是时间，纵轴是沿线的车站，每一列火车是图上一条斜线，斜率就是车速。两列车的线如果相交，意味着同一时刻出现在同一位置——也就是撞车。调度员的全部工作，用图像的语言说，就是让几百条线彼此永不相交。同方向快车追慢车，在图上是一条陡线逼近一条缓线，调度员要提前把慢车“掰”进侧线停一会——图上就是给缓线添一段水平线，等陡线先过去。你坐车时听到的“临时停车，待避列车”，背后都是这张图。

再来看一组带真实数字的估算，看看“面积”多么好使。高铁从静止加速到时速三百五十公里，大约需要五分钟。把时速三百五十公里换算成米每秒，约 \(97\) 米每秒。加速过程的速度—时间图近似一个三角形：底边是时间 \(300\) 秒，高是 \(97\) 米每秒。加速段驶过的距离就是这块三角形的面积：
\[
\Delta x\approx\frac{1}{2}\times 97\ \text{m/s}\times 300\ \text{s}\approx 1.5\times10^{4}\ \text{m},
\]
也就是约十五公里。不解任何方程，只数一块面积，我们就知道了：一列高铁从启动到全速，已经在不知不觉中开出了十几公里。这就是为什么高铁线路需要那么长的平直路段。

\begin{mathbridge}
如果你已经学过一点函数记号，可以把上面的极限写得更紧凑。瞬时速度写作
\[
v(t)=\lim_{\Delta t\to 0}\frac{x(t+\Delta t)-x(t)}{\Delta t}=\frac{dx}{dt},
\]
读作“\(x\) 对 \(t\) 的导数”。它就是一个函数：每个时刻对应一个斜率值。看个具体例子：设 \(x(t)=2t^{2}\)（位置以米为单位，时间以秒为单位），则
\[
\frac{x(t+\Delta t)-x(t)}{\Delta t}=\frac{2(t+\Delta t)^{2}-2t^{2}}{\Delta t}=4t+2\,\Delta t.
\]
当 \(\Delta t\) 趋于零，\(2\,\Delta t\) 这一项消失，得到 \(v(t)=4t\)。检查量纲（第 2 章的工具）：\(2t^{2}\) 的系数 \(2\) 其实带着“米每二次方秒”的单位，所以 \(4t\) 的单位是米每秒，确实是速度。

反过来，面积在极限下写成积分：
\[
\Delta x=\int_{t_1}^{t_2} v(t)\,dt.
\]
而“从位置求导得速度、对速度积分回位置”这对互逆关系，有一个响亮的名字——微积分基本定理。用本章的话说就是：\textbf{斜率图的面积，等于原图的高度差}。整条微积分大厦，地基就是这句图像语言。
\end{mathbridge}

\begin{challenge}
把“面积”用在速度—时间图上，还能白捡一个著名公式。匀加速运动的速度图是过原点的直线 \(v=at\)，到时刻 \(t\) 为止图下面积是三角形：
\[
\Delta x=\frac{1}{2}\times t\times at=\frac{1}{2}at^{2}.
\]
初中学过的匀加速位移公式，原来只是一块三角形面积，根本不用背。

再想深一层：斜率操作可以反复做。位置图求斜率得速度图，速度图再求斜率得到“速度变化有多快”的图——这个量在下一章正式登场，叫加速度。反过来，从加速度图量面积得速度，再量一次面积得位置。你的手机正是这样工作的：里面的加速度计每秒钟测量几百次加速度，芯片连续做两次“面积累加”，推算你的速度和位移。但工程上有个残酷的细节：每次测量都带一点小误差，而面积累加会把误差越滚越大——积分一次，误差线性增长；积分两次，误差按平方增长。所以纯靠加速度计推算位置，几分钟后就会漂得离谱，手机必须不断用卫星定位、无线信号来“校零”。斜率与面积的互逆，在数学里是精确的，在布满噪声的真实世界里却不对称：求斜率会放大噪声，求面积会积累偏差。这是图像语言走进工程时必上的一课。
\end{challenge}

\step{模型的边界}

图像语言极其强大，但正因为它强大，误用也极其普遍。下面四条边界，每一条都对应一类高频错误。

边界一：\textbf{图像是账，不是地图}。位置—时间图描述的是“什么时候在哪里”，不是“走了一条什么路”。沿直线来回走的物体，图线可以是起伏的曲线；而真正沿曲线运动的物体，只要沿路位置均匀变化，位置—时间图照样是一条直线。把图线的形状当成运动的轨迹，是初学者最常见的误读。图线弯曲，意思是变化的快慢在变，跟“拐弯”毫无关系——方向问题要等第 4 章引入向量才能说清楚。

边界二：\textbf{“看起来陡”不等于斜率大}。同一组数据，把横轴拉长，图就变平缓；把纵轴压缩，图就变陡峭。报纸广告里那些“业绩猛增”“暴跌”的图，很多是靠截断纵轴、拉伸横轴做出来的视觉效果。物理图像上唯一可信的是带着单位和数字的斜率值，不是眼睛的观感。看到任何一张图，先看两根轴各是什么量、什么单位、从零开始还是从中间截断。

边界三：\textbf{面积是带符号的账}。速度—时间图在横轴下方的部分，面积要记负。所以“面积”算出来的是位置的净变化（位移），不是总共走了多远（路程）。一个人向前走十米再走回十米，位移是零，路程是二十米——速度图上的正负面积抵消了。如果你想知道路程，得先把速度取绝对值再算面积。这组区别在下一章描述运动时会是主角。

边界四：\textbf{两点之间的连线是一种声明}。实验给你的永远是离散的点，把它们连成光滑曲线，是在声明“我相信中间时刻的行为也平滑地介于其间”。多数情况下这个声明是对的，但它不是白来的。水沸腾的那一分钟里，水温—时间图是一条平线，而在冰熔化的过程中温度—时间图同样出现平台——如果测量点太稀，你可能完全错过平台，把相变画成缓坡。采样多密才够，取决于你关心的过程变化多快，这是实验设计的一部分，不是画图软件替你决定的。

\step{回头解释开场现象}

现在回到那场百米决赛。把甲、乙两人的位置—时间图画在同一张图上。

\begin{center}
\begin{tikzpicture}[>=Stealth,scale=1.05]
  \draw[->] (0,0) -- (5.8,0) node[below left] {\(t\)（时间）};
  \draw[->] (0,0) -- (0,3.6) node[below left] {\(x\)（位置）};
  \draw[dashed] (0,3) node[left] {\(100\) 米} -- (5.2,3);
  \draw[thick] (0,0) -- (3.4,3) node[pos=0.75,above left] {甲};
  \draw[thick] (0,0) .. controls (1.6,0.5) and (2.8,2.5) .. (4.4,3) node[pos=0.8,below right] {乙};
  \draw[dashed] (3.4,3) -- (3.4,0) node[below] {\(t_{\text{甲}}\)};
  \draw[dashed] (4.4,3) -- (4.4,0) node[below] {\(t_{\text{乙}}\)};
\end{tikzpicture}
\end{center}

甲的图线几乎是一条陡的直线：起跑快，全程节奏均匀。乙的图线开头平缓——起跑反应慢，斜率小；中段越来越陡，斜率超过了甲，这就是雷达说的“乙中途速度更快”；两条线先后到达一百米高度，甲先，乙后。比赛的所有信息——谁起跑快、谁后程猛、谁先撞线——在这张图上一览无余。

开场的矛盾也随之化解。速度—时间图量的是“此刻变得多快”，对应乙的优势；而比赛规则在乎的是“谁先积累够一百米”，也就是谁的位置先变化一百米——这是面积的语言。乙的问题是：他速度更高的那段时间不够长，面积追回来之前，比赛已经结束了。用一句话模型的语言说：\textbf{乙赢在了斜率上，甲赢在了面积上}。“谁更快”之所以听起来有歧义，是因为日常语言把斜率和面积混叫成一个“快”字。图像不说话则已，一说话就把两种“快”分到了两根轴上。

这正是本章开头的承诺：遇到“一句话两种意思”的局面，不要吵，画图。

再用一组真实数据验收一下我们的工具。二〇〇九年柏林世锦赛，博尔特以九秒五八跑完一百米。整场比赛的“面积账”是：位置变化一百米。他的平均速度约为 \(100\div 9.58\approx 10.4\) 米每秒；而分段计时显示，他在六十米到八十米之间达到约每秒十二米的极速。也就是说，他的速度—时间图先从零快速爬升，在高位缓慢波动，图下总面积恰好一百米。假如有人全程保持平均速度十米每秒匀速跑，他的速度图是一条水平线，面积是矩形，跑完一百米需要十秒整——比博尔特慢。快出来的那零点四二秒，就藏在图线形状的差异里：博尔特把“面积”更早地攒够了。百米比赛比的不是谁速度数字大，而是谁更快积累够一百米的位移——这句话现在有了精确的图像含义。

\step{脑洞实验与挑战题}

\begin{thought}[你的人生是一条线]
把本章的工具推到一个极端尺度。想象一张巨大的图：横轴是时间，纵轴是你在一座城市里的位置。从今天往回画，你的每一次出门、上学、回家，都是图线上的一段。你会发现一个有点奇怪的视角：在这张图上，你从来没有“消失”过，你的人生是一条连续不断的线，从图的左边一直延伸到右边。就算你整天躺在床上不动，你的线也没有停——它只是变成一条水平线，安静地向时间深处延伸。物理学家给这条线起了个名字：世界线。

现在做两个极端推演。第一，什么样的世界线是“不可能”的？一条往左回头的线——同一个时刻出现两个你，那意味着回到过去。第二，一条完全竖直的线意味着什么？位置在变而时间完全不走，也就是不花任何时间就移动了距离，斜率无穷大。现实中所有物体的世界线都被限制在一个“陡度”以内，而这个限制不是工程问题，是宇宙的基本规则：存在一种任何世界线都不能超过的“最陡斜率”，它对应的速度就是光速。等你读到本书讲相对论的章节，会发现那里最重要的一张图——时空图——其实就是本章位置—时间图把两根轴的角色重新安排了一下，而“因果”“同时”“时间膨胀”这些大词，全都变成了图上斜率和角度的几何问题。你在本章学会的读图本领，到时候会直接派上用场。
\end{thought}

\begin{puzzles}
\item 电梯从一楼上升，中途在四楼开门停了十秒，然后继续到八楼；随后空驶回一楼。请定性画出它的位置—时间图，并回答：下行那一段的斜率与上行段斜率是什么关系？\textbf{思路提示}：停留时位置不变，图线是什么形状？下行时位置的变化方向反了，斜率的符号应该怎样？下行的绝对值看起来比上行大还是小，取决于什么？

\item 滴漏实验中，把直桶玻璃杯换成一个上粗下细的杯子（比如普通的锥形水杯），水流大小仍保持恒定。水面高度—时间图还是直线吗？如果不是，它向哪边弯？\textbf{思路提示}：斜率的意思是“水面升高的快慢”，而水流恒定意味着“单位时间增加的体积”恒定。杯子越往上越粗，同样的体积会把水面抬得更高还是更低？注意区分“体积的积累”和“高度的变化”——这是本章“面积与斜率不是一回事”的又一次现身。

\item 甲、乙两车在同一条直路上行驶，它们的速度—时间图都是直线：甲的图线从每小时六十公里水平延伸，乙的图线从零开始稳定上升，在某一时刻与甲的图线相交。交点意味着乙追上甲了吗？\textbf{思路提示}：交点是两条图线“高度相同”，高度在这张图上代表什么量？追上与否问的却是另一个量——两车各自图线下的面积差。交点那一刻，谁的面积更大？

\item 一颗弹性球从高处落下，落地反弹，再落再弹，每次弹起高度变低。定性画出它的速度—时间图（取向上为正），特别注意触地瞬间图线的样子。\textbf{思路提示}：下落时速度向下，为负，且越来越负；反弹后速度向上，为正。触地过程极短，速度从负猛跳到正，图线几乎是一个竖直的跳变——斜率在那瞬间极大。现实中跳变越陡，说明触地时间越短、速度改变越急，这个“急”的程度，正是后面章节里“力”的用武之地。
\end{puzzles}

这几道题没有一个需要你算出具体数字，但每一道都在考同一件事：你能不能让图替你思考。斜率回答“变得多快”，面积回答“积累多少”，符号报告方向——带着这三件兵器，下一章我们就要给“变化”本身装上方向，进入向量的世界。
